\documentclass[11pt]{article}
\usepackage{fontspec}
\usepackage[french,english]{babel}
\usepackage{amsmath}
\usepackage{amssymb}
\usepackage{amsthm}
\usepackage{mathrsfs}
\usepackage{tikz-cd}
\usetikzlibrary{babel}
\usepackage{enumitem}

\newcommand{\printedpage}[1]{\marginpar{\footnotesize p.~#1}}
\newcommand{\msp}[1]{\marginpar{\footnotesize [#1]}}

\newcommand{\Cat}{\mathrm{Cat}}
\newcommand{\HOT}{\mathrm{HOT}}
\newcommand{\hot}{\mathrm{hot}}
\newcommand{\pr}{\mathrm{pr}}
\newcommand{\id}{\mathrm{id}}
\newcommand{\Ob}{\mathrm{Ob}\,}
\newcommand{\Hom}{\mathrm{Hom}}
\newcommand{\uHom}{\underline{\mathrm{Hom}}}
\newcommand{\Liss}{\underline{\mathrm{Liss}}\,}
\newcommand{\cA}{\mathcal{A}}
\newcommand{\dd}{\mathrm{d}}
\newcommand{\lc}{\mathrm{lc}}
\newcommand{\adj}[2]{\mathrel{\underset{#2}{\overset{#1}{\rightleftarrows}}}}

\begin{document}

\printedpage{89}
% head: A. Grothendieck : Dérivateurs, Ch. VI, Édition des Universités de Montpellier II et Paris VII

in $(\Cat/S')(W^{\dd}_{S'})^{-1}$. Let us try to prove that this is commutative in $\Cat/S'$, not just in the category localized by $W^{\dd}_{S'}$. We can look at it as a diagram in $\Cat/S$ and prove its commutativity as such. This commutativity relation $vu = w$ then splits into two, according to the two factors $\Phi_S(X)$ and $S'$ of the target object $\Phi_S(X) \times_S S'$. For the factor $S'$, the relation simply means that $\alpha_{X'}$ is an $S'$-morphism, that's OK. For the factor $\Phi_S(X)$, we extend the diagram $(*)$ to the right by the square
\[
\begin{tikzcd}
X' \arrow[r, dotted, "\alpha_{X'}"] \arrow[dr, dotted] & \Phi_S(X') \times_S S' \arrow[r, "\pr_1^{\Phi_S(X')}"] \arrow[d, "\Phi_S(\pr_1 X)\times_S S'"] & \Phi_S(X') \arrow[d, "\Phi_S(\pr_1^X)"] \\
 & \Phi_S(X) \times_S S' \arrow[r, "\pr_1^{\Phi_S(X)}"] & \Phi_S(X),
\end{tikzcd}
\] % [sic] \Phi_S(\pr_1 X)\times_S S' -> \Phi_S(\pr_1^X)\times_S S', the X is printed on the line instead of as superscript of \pr_1
\msp{page 146}
and taking into account that the composite of the first horizontal row is $i_{X'}$, by construction of $\alpha_{X'}$, we are reduced to the commutativity of
\[
\begin{tikzcd}
X \times_S S' = X' \arrow[r, "i_{X'}"] \arrow[d, "i_X\times_S S'"'] & \Phi_S(X') \arrow[d, "\Phi_S(\pr_1^X)"] \\
\Phi_S(X) \times_S S' \arrow[r, "\pr_1^{\Phi_S(X)}"] & \Phi_S(X).
\end{tikzcd}
\tag*{$(*)$}
\]

Now the morphism $i_X : X \longrightarrow \Phi_S(X)$ in $\Cat/S$ gives rise to the commutative square (as does any arrow $i : X \longrightarrow Y$)
\[
\begin{tikzcd}
X \times_S S' \arrow[r, "\pr_1^X"] \arrow[d, "i_X\times_S S'"'] & X \arrow[d, "i_X"] \\
\Phi_S(X) \times_S S' \arrow[r, "\pr_1^{\Phi_S(X)}"] & \Phi_S(X),
\end{tikzcd}
\]
so that the lower composite of $(*)$ is none other than that of $(**)$, hence equal to the upper composite of the latter, and the commutativity of $(*)$ comes down to that of
\[
\begin{tikzcd}
X \times_S S' = X' \arrow[r, "i_{X'}"] \arrow[d, "\pr_1^X"'] & \Phi_S(X') \arrow[d, "\Phi_S(\pr_1^X)"] \\
X \arrow[r, "i_X"] & \Phi_S(X),
\end{tikzcd}
\]
which is obvious by functoriality of $Z \xrightarrow{i_Z} \Phi_S(Z)$, applied to the arrow $\pr_1^X : X' \longrightarrow X$ of $\Cat/S$.

% foot: Version 9 novembre 2001

\printedpage{90}
% head: A. Grothendieck : Dérivateurs, Ch. VI, Édition des Universités de Montpellier II et Paris VII

\msp{page 147}
We have thus proved that the functors
\[
\HOT_W(S') \adj{f_!}{f^*} \HOT_W(S)
\tag*{$(*)$}
\]
are adjoint to each other, with adjunction morphisms
\[
\left\{
\begin{array}{lcl}
\id & \xrightarrow{\ \alpha\ } & f^*f_! \\[1ex]
\id & \xleftarrow{\ \beta\ } & f_!f^*
\end{array}
\right.
\]
specified above. I presume that the same proof, working with the coarse localizers $W_S$, $W_{S'}$, instead of $W^{\dd}_S$, $W^{\dd}_{S'}$, would give that the corresponding functors
\[
\HOT^{\lc}_W(S') \adj{f_!^{\lc}}{f^*_{\lc}} \HOT^{\lc}_W(S)
\tag*{$(**)$}
\]
are adjoint to each other, with essentially ``the same'' adjunction morphisms $\alpha$, $\beta$. I do not feel like going over the preceding arguments again and prefer to \textit{deduce} it from the adjunction in the situation $(*)$. This will give me the opportunity to better understand the essential formal features of the situation. I shall put for simplicity
\msp{page 148}
\[
\begin{array}{rcl@{\qquad\qquad}rcl}
\HOT_W(S) & = & \cA, & \HOT_W(S') & = & \cA', \\
\HOT^{\lc}_W(S) & = & \cA_0, & \HOT^{\lc}_W(S') & = & \cA'_0,
\end{array}
\]
and I use the pairs of adjoint functors
\[
\cA \adj{\rho}{i} \cA_0, \qquad\qquad \cA' \adj{\rho'}{i'} \cA'_0,
\]
where $\rho$, $\rho'$ are localization functors, i.e. $i$, $i'$ fully faithful functors. I have a pair of adjoint functors
\[
\cA' \adj{f_!}{f^*} \cA,
\]
and I assume moreover

\begin{enumerate}
\item[a)] $f_!$ is compatible with the localizers of $\cA'$, $\cA$ defined by $\rho'$, $\rho$ respectively, hence defines a functor (well determined up to unique isomorphism)
\[
f_!^0 : \cA'_0 \longrightarrow \cA_0,
\]
thus giving rise to a commutative diagram (up to functorial isomorphism)

% foot: Version 9 novembre 2001

\printedpage{91}
% head: A. Grothendieck : Dérivateurs, Ch. VI, Édition des Universités de Montpellier II et Paris VII
\[
\begin{tikzcd}
\cA' \arrow[r, "\rho'"] \arrow[d, "f_!"'] & \cA'_0 \arrow[d, "f_!^0"] \\
\cA \arrow[r] & \cA_0.
\end{tikzcd}
\tag*{(I)}
\]
\texttt{[rather} $\cA \xrightarrow{\rho} \cA_0$\texttt{]}

\item[b)] ``Dually'', $f^* : \cA \longrightarrow \cA'$ is compatible with the full subcategories, essential images of the fully faithful functors $i$, $i'$, and then induces
\msp{page 149}
a functor, defined up to unique isomorphism:
\[
f_0^* : \cA_0 \longrightarrow \cA'_0,
\]
giving rise to a square of functors, commutative up to canonical functorial isomorphism
\[
\begin{tikzcd}
\cA_0 \arrow[r, "i"] \arrow[d, "f_0^*"'] & \cA \arrow[d, "f^*"] \\
\cA'_0 \arrow[r, "i'"] & \cA'.
\end{tikzcd}
\tag*{(II)}
\]
\end{enumerate}

\textit{Under these conditions, the pair of functors}
\[
\cA_0 \adj{f_!^0}{f_0^*} \cA'_0, \qquad \text{\texttt{[rather} } \cA'_0 \adj{f_!^0}{f_0^*} \cA_0, \text{\texttt{]}}
\]
\textit{is a pair of adjoint functors.} (NB when I speak of pairs of adjoint functors, the one written above the other is the left adjoint.)

\textsc{Proof}. Let
\[
\eta_0 \in \Ob \cA_0, \qquad\qquad \xi'_0 \in \Ob \cA'_0,
\]
we must define a bifunctorial bijection
\[
\Hom_{\cA_0}(f_!^0(\xi'_0), \eta_0) \simeq \Hom_{\cA'_0}(\xi'_0, f_0^*(\eta_0))
\]
For this, we are going to write $\xi'_0$ in the form
\[
\xi'_0 = \rho'(\xi')\,, \qquad\qquad \xi' \in \Ob \cA',
\]
it will suffice to establish the bifunctorial bijection by replacing $\xi'_0$ by $\rho'(\xi')$, i.e.
\[
\Hom_{\cA_0}(f_!^0(\rho'(\xi')), \eta_0) \simeq \Hom_{\cA'_0}(\rho'(\xi'), f_0^*(\eta_0))
\]

% foot: Version 9 novembre 2001

\printedpage{92}
% head: A. Grothendieck : Dérivateurs, Ch. VI, Édition des Universités de Montpellier II et Paris VII

\msp{page 150}
(Indeed, the adjunction formula is interpreted as an isomorphism of bifunctors, or, what amounts to the same thing, of functors $\cA'_0 \mathrel{\underset{\psi}{\overset{\varphi}{\rightrightarrows}}} \cA_0^{\wedge}$, now
\[
\begin{array}{ccc}
\uHom(\cA'_0, \cA_0^{\wedge}) & \longrightarrow & \uHom(\cA', \cA_0^{\wedge}) \\[1ex]
\varphi & \longmapsto & \varphi \circ \rho'
\end{array}
\]
is fully faithful, since $\rho'$ is a localization functor, so it amounts to the same thing to give oneself an isomorphism $\varphi \simeq \psi$, or $\varphi \circ \rho' \simeq \psi \circ \rho'$.) We must therefore transform the expression
\[
\begin{array}{cccc}
 & \Hom_{\cA_0}(\underbrace{f_!^0(\rho'(\xi'))}_{\rho f_!(\xi')}, \eta_0) & \underbrace{\simeq}_{\text{\footnotesize commutativity of (I)}} & \Hom_{\cA_0}(\rho f_!(\xi'), \eta_0) \\[3ex]
\underbrace{\simeq}_{\text{\footnotesize adjunction $\rho$, $i$}} & \Hom_{\cA}(f_!(\xi'), i(\eta_0)) & \underbrace{\simeq}_{\text{\footnotesize adjunction $f_!$, $f^*$}} & \Hom_{\cA'}(\xi', \underbrace{f^*i(\eta_0)}_{i'(f_0^*(\eta_0))}) \\[3ex]
\underbrace{\simeq}_{\text{\footnotesize commutativity of (II)}} & \Hom_{\cA'}(\xi', i'(f_0^*(\eta_0))) & \underbrace{\simeq}_{\text{\footnotesize adjunction $\rho'$, $i'$}} & \Hom_{\cA'_0}(\underbrace{\rho'(\xi')}_{\xi'_0}, f_0^*(\eta_0)),
\end{array}
\]
which is the announced bijection.

I summarize the essentials of the results obtained on $f_!$, $f_!^{\lc}$:

\textbf{Theorem 2.} \textit{Let $f : S' \longrightarrow S$ be an arrow in} $\Cat$\textit{, let us consider the functors}
\[
\begin{aligned}
\HOT_W(S') &\adj{f_!}{f^*} \HOT_W(S) \\
\HOT^{\lc}_W(S') &\adj{f_!^{\lc}}{f^*_{\lc}} \HOT^{\lc}_W(S),
\end{aligned}
\]
\msp{page 151}
\textit{where the functors $f_!$, $f_!^{\lc}$ are described in theorem} 1 \textit{(p. 136). The adjunction morphisms}
\[
\left\{
\begin{array}{lcl}
\id & \longrightarrow & f^*f_! \\[1ex]
\id & \longleftarrow & f_!f^*
\end{array}
\right.
\]
\textit{are described on pp. 138 to 140.}

We would like to make explicit the nice formula, obtained in passing
\[
\boxed{\;f_!f^*(\xi) = \hot_{S,W}(S'/S) \times_S \xi\;}\ ,
\]
for this we must make the right-hand side explicit. Let us note the

\textbf{Proposition 4.} \textit{There exists a unique functor}
\[
(\xi, \eta) \longmapsto \xi \times_S \eta : \HOT_W(S) \times \HOT_W(S) \longrightarrow \HOT_W(S)
\]

% foot: Version 9 novembre 2001

\printedpage{93}
% head: A. Grothendieck : Dérivateurs, Ch. VI, Édition des Universités de Montpellier II et Paris VII

\textit{with the property that for $X$, $Y$ varying in} $\Liss S$\textit{, we have an isomorphism bifunctorial in $X$, $Y$}
\[
\hot_{S,W}(X) \times_S \hot_{S,W}(Y) \simeq \hot_{S,W}(X \times_S Y)
\tag*{$(*)$}
\]
(\footnote[72]{\textbf{NB} The multiplication $\times_S$ in $\HOT_W(S)$ is associative, commutative, unital \textit{etc., etc.}, since this is the case in $\Liss S$.}).

This is formal, using only that on $\Liss S$, the bifunctor
\[
(X, Y) \longrightarrow X \times_S Y
\]
is compatible with the localizers $W^{\dd}_S \times W^{\dd}_S$ and $W^{\dd}_S$. To see this, we are reduced to seeing that if $X \xrightarrow{u} X'$ in $\Liss S$ is in $W^{\dd}_S$, then the same holds for $X \times_S Y \xrightarrow{\ u\times_S \id_Y\ } X' \times_S Y$, for every
\msp{page 152}
$Y$ in $\Liss S$. In fact, we use only the smoothness of $X$, $X'$, not that of $Y$, and this is immediate, for if $X \longrightarrow X'$ is in $W^{\dd}_S$, $X \times_S Y \longrightarrow X' \times_S Y$ is in $W^{\dd}_Y$, hence also in $W^{\dd}_S$ (cf. theorem 1 applied to $Y \longrightarrow S$). This remark suggests at the same time:

\textbf{Corollary 1. (False)} \textit{The bifunctorial formula} $(*)$ \textit{remains valid for $X$ in} $\Liss S$\textit{, $Y$ arbitrary in} $\Cat/S$\textit{, and also in the reverse case.} (\footnote[73]{This statement is surely \textit{false!} Counterexample p. 153, cf. corollary 1 corrected on page 154-156.})

In the first case, let us consider
\[
Y \xrightarrow{i_Y} \Phi(Y)
\]
in $W^{\dd}_S$, hence an isomorphism in $(\Cat/S)(W^{\dd}_S)^{-1}$, we shall therefore have
\[
\begin{array}{lcl}
\hot_{S,W}(X) \times_S\ \hot_{S,W}(Y) & \xrightarrow{\ \sim\ } & \hot_{S,W}(X) \times_S\ \hot_{S,W}(\Phi(Y)) \\[1ex]
 & \underbrace{\xrightarrow{\ \sim\ }}_{\text{\footnotesize since $X$ \textit{and} $\Phi(Y)$ smooth over $S$}} & \hot_{S,W}(X \times_S \Phi(Y)),
\end{array}
\]
now I say that
\[
X \times_S Y \xrightarrow{\ X\times_S i_Y\ } X \times_S \Phi(Y)
\]
is again in $W^{\dd}_S$. It suffices to see that it is in $W^{\dd}_X$ (cf. theorem 1), and for this, that the smooth base change $X \longrightarrow S$ induces $\Cat/S \longrightarrow \Cat/X$ compatible with the localizers $W^{\dd}_S$, $W^{\dd}_X$. But this is \textit{false} (take $X = \{s\}$).

\msp{page 153}
Let us show that the corollary is false. I am going to take $Y = \{s\}$, the one-point category at $s$, then $\Phi(Y)$ is the fibered category with discrete fibers defined by the representable functor $F$ defined by $s$, so it is $S/F$ and, in this case $S/s$: % [sic] les corollaire -> le corollaire, agreement; occurence -> occurrence, spelling
\[
\Phi(\{s\}) \simeq S/s \qquad\quad \text{as a category over } S.
\]

% foot: Version 9 novembre 2001

\end{document}
